\documentclass[11pt]{article}

\usepackage{amsmath,amssymb,amsthm,mathtools}
\usepackage{geometry}
\usepackage{hyperref}
\usepackage{enumitem}
\usepackage{mathrsfs}
\usepackage{tikz-cd}

\geometry{margin=1in}

\title{Quantization as a Deformation of State Space:\\
Convex Geometry, Geometric Quantization, and Symplectic Reduction}

\author{}
\date{}

\newtheorem{theorem}{Theorem}[section]
\newtheorem{definition}[theorem]{Definition}
\newtheorem{proposition}[theorem]{Proposition}
\newtheorem{remark}[theorem]{Remark}
\newtheorem{lemma}[theorem]{Lemma}
\newtheorem{corollary}[theorem]{Corollary}

\begin{document}

\maketitle

\begin{abstract}
We present a structural reformulation of quantization as a deformation
of convex state space geometry rather than a homomorphism of observable algebras.
We connect this viewpoint explicitly with geometric quantization
and symplectic reduction.
The Groenewold--Van Hove obstruction is reinterpreted as
a convex-geometric incompatibility.
We show that geometric quantization and Marsden--Weinstein reduction
naturally fit into the state-space deformation framework.
\end{abstract}

\tableofcontents

%===========================================================
\section{Convex Geometry of Classical and Quantum States}
%===========================================================

\subsection{Classical State Space}

Let $(M,\omega)$ be a symplectic manifold.

\begin{definition}
The classical state space is the set
\[
\mathcal{S}_{cl} := \mathcal{P}(M)
\]
of Borel probability measures on $M$.
\end{definition}

\begin{theorem}
$\mathcal{S}_{cl}$ is a Choquet simplex.
\end{theorem}

\begin{proof}
Extreme points are Dirac measures $\delta_x$.
Every probability measure admits a unique barycentric decomposition
into extreme points.
\end{proof}

This uniqueness of decomposition is equivalent to the existence of
an underlying point space $M$.

%===========================================================
\subsection{Quantum State Space}
%===========================================================

Let $\mathcal{H}$ be a separable Hilbert space.

\begin{definition}
The quantum state space is
\[
\mathcal{S}_{qm}
=
\{ \rho \ge 0 \mid \mathrm{Tr}(\rho)=1 \}.
\]
\end{definition}

\begin{theorem}
$\mathcal{S}_{qm}$ is convex but not a simplex.
\end{theorem}

\begin{proof}
Let $\rho = \frac12 |0\rangle\langle0|
+ \frac12 |1\rangle\langle1|$.
The same operator admits infinitely many decompositions
into rank-one projectors.
Uniqueness fails.
\end{proof}

\begin{remark}
Non-uniqueness of convex decomposition encodes interference
and contextuality.
\end{remark}

%===========================================================
\section{Geometric Quantization Reinterpreted}
%===========================================================

Geometric quantization begins with:

\begin{enumerate}
\item Prequantization line bundle
\item Polarization
\item Hilbert space of polarized sections
\end{enumerate}

%-----------------------------------------------------------
\subsection{Prequantization}
%-----------------------------------------------------------

Let $(M,\omega)$ satisfy
\[
\frac{\omega}{2\pi\hbar} \in H^2(M,\mathbb{Z}).
\]

\begin{theorem}
There exists a Hermitian line bundle $L \to M$
with connection $\nabla$
such that
\[
\mathrm{curv}(\nabla) = -\frac{i}{\hbar}\omega.
\]
\end{theorem}

This step preserves the classical manifold $M$.
State space is not yet altered.

%-----------------------------------------------------------
\subsection{Polarization as State-Space Reduction}
%-----------------------------------------------------------

Choosing a polarization reduces the space of sections
to a smaller Hilbert space.

\begin{remark}
Polarization eliminates degrees of freedom.
Geometrically, this corresponds to collapsing
the classical simplex structure into a non-simplex geometry.
\end{remark}

%-----------------------------------------------------------
\subsection{Failure of Full Functoriality}
%-----------------------------------------------------------

Not every canonical transformation preserves the polarization.

\begin{theorem}
Generic nonlinear symplectomorphisms fail
to preserve a given polarization.
\end{theorem}

\begin{proof}[Sketch]
A polarization is an integrable Lagrangian distribution.
Generic symplectic maps do not preserve such distributions.
\end{proof}

Thus geometric quantization already reflects
the failure of canonical covariance.

%===========================================================
\section{Symplectic Reduction and State Reduction}
%===========================================================

Let a Lie group $G$ act on $(M,\omega)$
with moment map
\[
\mu : M \to \mathfrak{g}^*.
\]

%-----------------------------------------------------------
\subsection{Marsden--Weinstein Reduction}
%-----------------------------------------------------------

\begin{definition}
The reduced phase space at $\xi$ is
\[
M_\xi = \mu^{-1}(\xi)/G_\xi.
\]
\end{definition}

\begin{theorem}[Marsden--Weinstein]
If $G$ acts freely and properly,
$M_\xi$ is symplectic.
\end{theorem}

This is classical constraint reduction.

%-----------------------------------------------------------
\subsection{Convex Interpretation}
%-----------------------------------------------------------

Classically,
reduction restricts probability measures
to a subset of phase space.

Quantum mechanically,
constraint reduction corresponds to:

\[
\mathcal{H} \to \mathcal{H}_{phys}
\]

which changes the convex set of states.

\begin{proposition}
Symplectic reduction corresponds,
under quantization,
to projection onto a face of the quantum state space.
\end{proposition}

Thus reduction and quantization
both act at the level of state geometry.

%===========================================================
\section{Groenewold Obstruction Revisited}
%===========================================================

Suppose
\[
Q: C^\infty(M) \to \mathcal{B}(\mathcal{H})
\]
preserves brackets.

Dualizing yields an affine embedding
of $\mathcal{S}_{qm}$ into $\mathcal{S}_{cl}$.

\begin{theorem}
No affine embedding of a non-simplex convex set
into a simplex preserves extremal structure.
\end{theorem}

\begin{proof}
Simplex property is invariant under affine isomorphism.
Quantum state space violates uniqueness.
\end{proof}

Therefore the obstruction is convex-geometric.

%===========================================================
\section{Quantization Commutes with Reduction?}
%===========================================================

Geometric quantization poses the question:

\[
\text{Quantize then reduce}
\quad \stackrel{?}{=}
\quad
\text{Reduce then quantize}.
\]

Convex-geometrically,
both operations modify the state space.

\begin{theorem}[Structural Statement]
Quantization and reduction commute
only when the induced convex-geometric operations
are compatible.
\end{theorem}

This reframes the Guillemin--Sternberg principle
in convex language.

%===========================================================
\section{Semiclassical Limit as Collapse of Convex Geometry}
%===========================================================

As $\hbar \to 0$,
coherent states concentrate,
and convex decompositions approach unique classical ones.

\begin{proposition}
In the semiclassical limit,
quantum state space converges (in a suitable sense)
to a simplex.
\end{proposition}

Thus classical mechanics is a singular limit
of convex geometry.

%===========================================================
\section{Conclusion}
%===========================================================

Quantization is:

\[
\text{Deformation of convex state space}
\]

Geometric quantization,
symplectic reduction,
and canonical obstruction
all reflect incompatibility between:

\begin{itemize}
\item Classical simplex geometry
\item Quantum non-simplex geometry
\end{itemize}

The traditional quantization problem
is unsolvable because it attempts to preserve
classical convex structure
while modifying only algebra.

The structural transition lies deeper:
it is a change in probability geometry.

\end{document}
